% Transcripción de texto y corrección de errores ortográficos

\section*{DUCTO PARA LA FUENTE RADIACTIVA}

\begin{figure}
    \centering
    \includegraphics[width=0.8\textwidth]{ducto.png} % Ajustar ruta y escala adecuadamente
\end{figure}

El ducto estándar tiene 7 m de longitud y sirve como guía para que por él se traslade la fuente desde el cabezal hasta el punto en el cual se va a tomar la radiografía. El ducto consiste de una manguera tejida, flexible, a prueba de intemperie y con accesorios de acero inoxidable en sus extremos.

Un extremo cuenta con una cuerda apropiada para poder enroscarse y fijarse al cabezal. El otro extremo cuenta con una tapa de cabeza de muñón que limita el recorrido de la fuente e impide que ésta salga. La tapa de cabeza de muñón es removible, lo que permite las siguientes operaciones:

\begin{enumerate}[label=\alph*)]
    \item Conectar en el extremo del ducto el colimador de tungsteno.

\item Conectar la tapa de cabeza de muñón
Conectar la tapa de cabeza de muñón directamente a la salida del cabezal para tomar radiografías panorámicas dentro de un cuarto de radiografía o detrás de una barrera blindada.

\end{enumerate}
